\documentclass[11pt]{article}
\usepackage{amsmath}
\usepackage{booktabs}
\usepackage{parskip}
\usepackage[margin=1in]{geometry}

\title{Restoring Backslashes}
\author{A. Author}
\date{October 2026}

\begin{document}
\maketitle
\tableofcontents

\section{Introduction}
\label{sec:intro}

\noindent Crystals diffracted to 1.5 \AA{} in 80\% of trials.
The structure is described in \cite{Smith2020} and in Section~\ref{sec:methods}.
We thank Smith \& Jones for \textbf{bold} and \emph{emphasized} help.

\subsection*{Background}

The energy is $E = \frac{1}{2} m v^2$ and the angle is $\theta \leq \pi$.
The sum $\sum_{i=1}^{n} \alpha_i \to \infty$ converges when $x_{in} \in A$.
The alpha particles and beta sheets in this prose are left alone.

\begin{equation}
  \label{eq:main}
  \Delta G = \Delta H - T \Delta S \quad \text{for all} \quad T > 0
\end{equation}

\begin{align}
  a &= \sqrt{b^2 + c^2} \\
  \nabla \cdot \vec{E} &= \frac{\rho}{\epsilon_0}
\end{align}

\section{Methods}
\label{sec:methods}

\begin{itemize}
  \item First item with an item inside the sentence.
  \item[b)] Second item.
\end{itemize}

\begin{table}[h]
  \centering
  \caption{Data collection statistics.}
  \begin{tabular}{lr}
    \toprule
    Parameter & Value \\
    \midrule
    Resolution & 1.5 \\
    Completeness & 99.1 \\
    \bottomrule
  \end{tabular}
\end{table}

\begin{figure}[h]
  \centering
  \includegraphics[width=0.8\textwidth]{figures/my_figure.png}
  \caption{A figure.}
\end{figure}

\begin{verbatim}
section{this stays as it is}
\end{verbatim}

See \url{https://example.org/a%20b} for details.

\bibliographystyle{plain}
\bibliography{refs}
\end{document}
